\begingroup
\NoSourceHead
\SourceLeaf{239}
\pdfbookmark[1]{Briefwechsel: Pfaff an Gauss, 24. November 1799}{x1-correspondence-1}
\begin{center}\Large\bfseries\textls[120]{BRIEFWECHSEL.}\end{center}
\NumberedItem{1}
\begin{center}\large\bfseries\textsc{Pfaff} an \textsc{Gauss}.\end{center}
\begin{flushright}\{Helmstedt, 24. November 1799.\end{flushright}
\begin{center}Theuerster Freund\end{center}
\noindent\hspace*{1.2em}.\quad .\quad .\quad .\quad .\quad .\quad .\quad .\quad .\quad .\quad .\quad .\quad .\quad .\quad .\quad .\quad .\quad .\quad .\quad .\quad .\quad .\quad .\quad .\quad .\par
\noindent Nun noch ein paar Worte, soweit es die Kürze der Zeit gestattet, von den mathematicis, welche Ihr Brief enthält.\quad .\quad .\quad .\quad .\quad .

\par\indent Sehr interessant war mir die besondere Bemerkung, welche Sie über das arithm. geom. Mittel zwischen $1$ und $\surd2$ gemacht haben, wobey die Übereinstimmung gewiss nicht zufällig und isolirt ist, sondern noch mehr in sich schliessen wird. — So leicht das arithm. harm. Mittel zu finden und zu beweisen ist, so wenig scheint bey der ersten Ansicht klar zu seyn, wie das a. g. M. zu bestimmen sey. Ich hätte gewünscht, Musse zu haben, mich auf diese curiose Speculation einzulassen. Ich bin aber in diesem Winter zu sehr mit meinen Vorlesungen beschäftigt, so dass ich kaum Zeit genug haben werde, meine Abhandlung über die meth. tang. invers. vollends zum Druck ins reine zu bringen, welche ich auch inzwischen wieder ganz liegen lassen musste. Indessen habe ich doch der Versuchung nicht widerstehen können, eben jezt, da mich Ihr Brief wieder daran erinnerte, wenigstens einen Versuch zu machen, wovon ich doch ein paar Worte sagen will. Ich sezte (wodurch die Untersuchung nicht beschränkt wird) $a=1$, $b=1+x$; so wird
\[
a'=\surd(1+x),\quad b'=1+\tfrac12x.
\]
\clearpage
\ArithmeticHead{233}{R}{PFAFF AN GAUSS.}
\SourceLeaf{240}
\noindent Nun suchte ich überhaupt die $a$ und $b$ durch Reihen nach $x$ auszudrücken, und das Gesez der Coefficienten zu bestimmen.

\par\indent Es sei
\[
\begin{aligned}
a^n&=1+\alpha^n x+\beta^n x^2+\gamma^n x^3+\delta^n x^4\ldots\phantom{,}\\
b^n&=1+A^n x+B^n x^2+C^n x^3+D^n x^4\ldots,
\end{aligned}
\]
\noindent so ergeben sich 2 Reihen von Gleichungen für die Coëfficienten: die eine ist leicht:
\[
A^{n+1}=\frac{A^n+\alpha^n}{2};\quad
B^{n+1}=\frac{B^n+\beta^n}{2};\quad
C^{n+1}=\frac{C^n+\gamma^n}{2};\quad\text{etc.}
\]
\noindent bey der anderen ist wegen des Infinitinomii die Rechnung verwickelter, ich habe sie nicht weiter fortführen können, als bis $\delta$ und $D$, bin aber zweifelhaft, ob, was ich bis dahin fand, richtig ist, da ich zu eilig und unordentlich rechnete. Ich erhielt nehmlich für \textls[120]{jedes} $n$ (exc[eptis] exc[ipiendis])
\[
\alpha^n=A^n=\tfrac12;\quad
\beta^n=B^n=-\tfrac1{16};\quad
\gamma^n=C^n=\tfrac1{32};\quad
\delta^n=D^n=-\frac5{2.128}\mathpunct{\raisebox{.5ex}{.}}
\]
\noindent Es wäre doch auffallend, dass von $a''$ und $b''$ an die Reihen wenigstens in den 5 ersten Gliedern übereinstimmen. Ich wäre neugierig zu wissen, wie es weiter geht, was sich für ein Gesez ergibt, ob etwa für $n=\infty$ die Reihe sich summiren lässt, wenigstens für gewisse Werthe von $x$. Aber die Zeit hat kaum zugereicht, so weit zu rechnen, und fast wäre mir die Zeit zu kurz geworden, um noch an Sie zu schreiben.

\par\indent Ich werde übrigens diese Aufgabe bey mehrerer Musse nicht aus dem Sinn verliehren, ob gleich diese Untersuchung, die weiter führen kann, bey Ihnen in guten Händen ist. — Erhalten Sie ferner Ihre freundschaftlichen, mir sehr schätzbaren Gesinnungen
\begin{flushright}
Ihrem ergebensten Fr[eunde]\\
\textsc{J. F. Pfaff.}\}
\end{flushright}
\clearpage
\endgroup
